\documentclass[11pt]{article}

\usepackage{amsmath,amssymb}
\usepackage{amsthm}
\usepackage[margin=1in]{geometry}
\usepackage{hyperref}
\usepackage[round,authoryear]{natbib}

\newtheorem{result}{Result}

\title{Pass-Through, Welfare, and Incidence under Product Returns\\
  in Perfect Competition}
\author{Draft notes}
\date{\today}

\begin{document}
\maketitle

%=============================================================================
\section{Introduction}
%=============================================================================

Who bears a tax, and how much of it passes through into the price, is the
classic question of tax incidence \citep{weyl2013pass}. A platform's commission
is one example. Amazon deducts a referral fee from each sale, a commission of 15
percent of the sale price in many categories, so the seller receives the price
net of the fee. The fee therefore works like a sales tax, and without returns
its pass-through and incidence are well studied in the literature. Online sales, however, are often
returned, with retailers expecting about one in five to come
back.\footnote{National Retail Federation and Happy Returns,
  \emph{2025 Retail Returns Landscape}, October 15, 2025,
  \url{https://nrf.com/research/2025-retail-returns-landscape}.
  The report estimates that 19.3 percent of online sales will be returned in
  2025.} On a kept sale the seller pays the whole commission. On a return the
seller refunds the price in full and gets most of the commission back, but
Amazon keeps part of it, at most 20 percent, or 3 percent of the price at a 15
percent fee.\footnote{Amazon, \emph{Selling on Amazon pricing},
  \url{https://sell.amazon.com/pricing}, and \emph{North America selling fees},
  \url{https://sell.amazon.com.sg/north-america/pricing}. Amazon calls the
  amount it keeps the refund administration fee, the lesser of five dollars or
  20 percent of the referral fee.} The seller also bears the cost of handling
the return. Returns thus change what the fee costs the seller. It is not clear
how that changes the pass-through and the burden borne by sellers and
consumers, or whether the platform should refund only part of the commission on
a return. We study this in a competitive market, taking the fee schedule as
given. Each buyer decides whether to buy and, after learning the fit, whether
to keep the good, and the fee reaches both decisions only through the price.

The standard case against a tax is that it opens a gap between what the buyer
will pay and what the good costs the seller, so some beneficial trades do not
happen. That argument has one margin, the decision to buy. Returns add a
second, the decision to keep the good. What the seller pays in fees also
depends on how many units come back. A higher price makes the buyers who remain more likely to return, because the
good must fit better to be worth keeping. It also changes who buys, and the
buyers it deters may return more often or less often than those who remain. So
returns can rise or fall with the price, and no-return formulas, which hold the
return rate fixed, do not carry over.

Counting the platform's revenue as part of social surplus, a small fee is only
a transfer when nothing is returned, so it leaves social surplus unchanged at
the margin. Returns break this benchmark. The fee changes the price, the price
changes how often the good comes back, and a return uses real resources in
handling while the refund itself is only a transfer. The share of the fee that
the platform keeps on a refund changes how much the price moves, not whether it
moves, since even a fee that is fully paid back on a return is still paid on
every unit the buyer keeps. Whether society gains or loses is therefore open,
and it is a different question from whether buyers and sellers lose.

The paper asks who bears a fee that is partly refunded on returns and whether
it is socially costly. It does so for a small change in the fee, in four steps.
First, how does the fee pass through into the price, and how do buyers and
sellers divide the change in their own surplus? A rise in production cost and a
rise in handling cost serve as benchmarks, since the first lowers the seller's
margin on every unit, the second only on returned units, and the fee lies
between the two. Second, counting the platform's revenue, does any loss to
society come from the decision to buy or from the decision to keep the good,
and in which direction does that decision go wrong? Third, what is missed by
ignoring returns and treating the fee as an ordinary tax? Fourth, if the
platform keeps a larger share of the fee on a refund, does society gain or
lose? The answers depend only on the return rate and on how demand, supply, and
returns respond to the price.

%=============================================================================
\section{Setting}
%=============================================================================

A homogeneous product sells at list price $p$. Consumers indexed by $i$
differ in return hassle $\kappa_i$, the disutility of sending the good back,
and in the distribution $G_i$ of post-purchase match quality $u$. After
buying, type $i$ draws $u\sim G_i$. The consumer keeps the item when
$u-p\geq -\kappa_i$, and then earns $u-p$. Otherwise the consumer returns
the item and earns $-\kappa_i$. The refund equals the purchase price.

Expected surplus from buying at price $p$ is
\[
  cs_i(p)=\mathbb{E}_{u\sim G_i}\bigl[\max(u-p,\,-\kappa_i)\bigr].
\]
Willingness to pay $v_i$ is defined by $cs_i(v_i)=0$. Type $i$ buys if and
only if $p\leq v_i$. Across types, $v_i$ has distribution $F$ with density
$f$, so gross demand is $D(p)=1-F(p)$.

The aggregate return rate $r(p)$ is the average return probability among
buyers at price $p$. Type $i$'s return probability is
\[
  \phi_i(p)=\Pr(u<p-\kappa_i).
\]
Buyers with the same willingness to pay $v$ can differ in $\kappa_i$ and
$G_i$, so they need not share a return probability. Their average return
probability is
\[
  \phi(v,p)\equiv\mathbb{E}\bigl[\phi_i(p)\bigm|v_i=v\bigr],
\]
which depends on $v$ when the joint distribution of $(\kappa_i,G_i)$ varies
with willingness to pay. Averaging over buyers at $p$ gives
\[
  r(p)=\frac{1}{D(p)}\int_p^\infty \phi(v,p)\,dF(v).
\]

Differentiating $r(p)$ and using $D'(p)=-f(p)$ splits the price response
into a behavioral term and a selection term,
\[
  r'(p)
  =
  \underbrace{
    \frac{1}{D(p)}\int_p^\infty \frac{\partial \phi(v,p)}{\partial p}\,dF(v)
  }_{\text{Behavioral channel}}
  +
  \underbrace{
    \frac{f(p)}{D(p)}\bigl(r(p)-\phi(p,p)\bigr)
  }_{\text{Selection channel}},
\]
where $\phi(p,p)$ is the average return probability among types at the
purchase margin $v_i=p$.

The behavioral term is the change in return probability among buyers who
remain in the market. A higher $p$ raises the keep threshold, so a return
becomes strictly more likely for those buyers.

The selection term is the change in who buys. A higher price drops buyers at
the margin $v_i=p$. For arbitrary $(\kappa_i,G_i)$ this term is unsigned,
because marginal buyers may return more or less often than the average
inframarginal buyer.

An additive match signs that term. Under $u=v_i+\varepsilon$, with
$\varepsilon$ and return hassle both independent of $v_i$, $\phi(v,p)$
decreases in $v$. Marginal buyers then return more often than inframarginal
buyers, so $\phi(p,p)\geq r(p)$ and selection is non-positive. Behavioral
and selection then move $r'(p)$ in opposite directions. Below, $r'(p)$ is
left unsigned.

Sellers face that return rate under perfect competition. Producer $j$ takes
$p$ as given, chooses output at cost $C_j(q_j)$, and pays a net handling
cost $h$ on each
returned unit. Net handling is gross handling cost minus scrap value. Net
handling may be negative. We maintain $p+h>0$, which holds whenever scrap
value does not exceed list price.

The list price is what the buyer pays the seller. The platform charges an
ad valorem fee on that receipt at rate $\tau$. When the good is kept, a
share $\tau$ of $p$ goes to the platform and the seller keeps $(1-\tau)p$.

When the good is returned, the seller refunds the full price, so the buyer's
net payment is zero. A share $\beta\in[0,1]$ of the fee is not returned to
the seller. The seller pays $\beta\tau p$ on that unit and pays the handling
cost $h$.\footnote{Amazon's refund
  administration fee is this rule for a referral fee.  On a return Amazon
  keeps the lesser of \$5 or 20 percent of the referral fee, and gives the
  rest back to the seller.  When the \$5 cap does not bind, $\beta=0.2$.}
None of the fee comes back to the seller when $\beta=1$. All of it comes
back when $\beta=0$.

The seller's payoff is $(1-\tau)p$ on a kept unit and $-\beta\tau p-h$ on a
returned unit. With return rate $r$, expected payoff per unit sold is
\begin{equation}
  (1-r)(1-\tau)p - r(\beta\tau p + h).
  \label{eq:seller-payoff}
\end{equation}

%=============================================================================
\section{Consumer side}
%=============================================================================

Consumer surplus at list price $p$ integrates expected surplus over buyers.
Write $cs(v,p)$ for expected surplus among types with willingness to pay
$v$, so $cs(v,v)=0$. Then
\begin{equation}
  CS(p)=\int_p^\infty cs(v,p)\,dF(v).
\end{equation}

The price derivative of consumer surplus equals minus kept demand.
Differentiating in $p$ gives an extensive-margin term, from types at the
purchase margin, and an intensive-margin term, from surplus among
inframarginal buyers,
\[
  CS'(p)
  = -cs(p,p)\,f(p)+\int_p^\infty cs_p(v,p)\,dF(v),
\]
where $cs_p(v,p)\equiv \partial cs(v,p)/\partial p$. The extensive margin is
zero because $cs(p,p)=0$. On the intensive margin, a higher price lowers a
buyer's surplus only when she keeps the good. Among buyers with willingness
to pay $v$,
\[
  cs_p(v,p)
  =
  -\Pr(\text{keep}\mid v_i=v).
\]
Aggregating over buyers gives
\[
  CS'(p)=\int_p^\infty cs_p(v,p)\,dF(v)
  = -D(p)\bigl(1-r(p)\bigr),
\]
by the definitions of $D(p)$ and $r(p)$.

A price path from $p_0$ to $p_1$ therefore changes consumer surplus by
\[
  \Delta CS = -\int_{p_0}^{p_1} D(p)\bigl(1-r(p)\bigr)\,dp.
\]
A kept unit loses the price increase. A returned unit loses nothing. Buyers
differ in how often they keep, so the same price increase cuts surplus by
different amounts for different buyers. Those differences average into
$r(p)$. $CS'(p)$ is then determined by $D(p)$ and $r(p)$, and $\Delta CS$ by the
paths of those two aggregates.

Integrating that derivative from $p$ up to a price $\bar{p}$ at which nobody
buys gives the level of consumer surplus,
\begin{equation}
  CS(p)=\int_p^{\bar{p}} D(z)\bigl(1-r(z)\bigr)\,dz.
\end{equation}

%=============================================================================
\section{Supply side}
%=============================================================================

Set the fee aside. A gross sale brings revenue $p$ if the consumer keeps the
item, and it costs the refund $p$ plus net handling $h$ if the item is
returned. Among gross sales at list price $p$, share $r(p)$ are returned, so
expected revenue per gross unit, the \emph{effective margin}, is
\begin{equation}
  \mu(p,h) \equiv p - r(p)(p+h).
  \label{eq:mu-def}
\end{equation}
Firm $j$ takes $p$ as given, so this margin enters the supply decision as a
fixed number $\mu$.

Given $\mu$, firm $j$ chooses output to solve
\[
  \max_{q_j \geq 0}\;\; \mu\,q_j - C_j(q_j),
  \qquad
  \text{FOC: } C_j'(q_j)=\mu \;\;(\text{if } q_j>0),
\]
with $C_j(0)=0$. Write $q_j(\mu)$ for firm $j$'s supply, equal to zero when
$\mu$ is below its minimum marginal cost. Strictly increasing marginal cost
implies $q_j'(\mu)>0$ for active firms.

Market supply and industry profit sum these choices.
\begin{equation}
  S(\mu) \equiv \sum_j q_j(\mu),
  \qquad
  PS(\mu) \equiv \sum_j \bigl[\mu\,q_j(\mu)-C_j\bigl(q_j(\mu)\bigr)\bigr].
  \label{eq:agg-def}
\end{equation}

Strictly increasing marginal cost implies $S'(\mu)>0$, so quantity supplied
rises with the effective margin. Market supply therefore falls in the list
price whenever $\mu(p,h)$ falls in $p$.

Industry profit is the area to the left of market supply. The envelope
theorem implies that firm $j$'s profit rises with $\mu$ at rate
$q_j(\mu)$. Summing over firms,
\[
  PS'(\mu)=\sum_j q_j(\mu)=S(\mu).
\]
With no production at $\mu=0$, that derivative integrates to
\begin{equation}
  PS(\mu)=\int_0^\mu S(z)\,dz.
  \label{eq:PS-integral}
\end{equation}
At equilibrium,
$PS=\int_0^{\mu(p,h)}S(z)\,dz$.

%=============================================================================
\section{A seller-remitted ad valorem tax}
%=============================================================================
\label{sec:sales-tax}

The expected payoff \eqref{eq:seller-payoff} is the seller's effective margin
once the fee is included,
\begin{equation}
  \mu^{s}(p,\tau,\beta,h)
  =
  (1-r(p))(1-\tau)p - r(p)(\beta\tau p + h).
  \label{eq:mu-s}
\end{equation}
The buyer's return decision depends only on $p$, so $r=r(p)$. The refund on
a return is the list price, and the fee is remitted by the seller. Rewriting
\eqref{eq:mu-s} in the pre-fee margin gives
\begin{equation}
  \mu^{s}
  =
  \mu(p,h) - \tau p\bigl[1-(1-\beta)r(p)\bigr].
  \label{eq:mu-s-split}
\end{equation}
The term in brackets is the share of the fee the seller expects to pay. That
share is the whole fee on a kept unit, and the share $\beta$ on a returned
unit.

The fee reaches consumers only by moving the equilibrium list price.
Market clearing is $D(p)=S(\mu^{s})$. Consumer surplus depends on the list
price. Producer surplus is the area under supply at $\mu^{s}$,
\[
  CS=CS(p),
  \qquad
  PS=\int_0^{\mu^{s}}S(z)\,dz.
\]

Tax derivatives are taken at $\tau=0$. At that rate, the fee term in
\eqref{eq:mu-s-split} has derivative zero with respect to $p$, so $\mu_p$ is
the pre-fee derivative of $\mu(p,h)$. Holding $p$ fixed, one percentage point
lowers the margin by the expected fee share times the list price,
\[
  \mu^{s}_{\tau}=-p\bigl[1-(1-\beta)r(p)\bigr].
\]
Write $Q\equiv D(p)$ and $\rho_{\tau}\equiv dp/d\tau$. Then
\begin{equation}
  \frac{dCS}{d\tau}
  =-Q\bigl(1-r\bigr)\rho_{\tau},
  \qquad
  \frac{dPS}{d\tau}
  =Q\Bigl(\mu_p\rho_{\tau}-p\bigl[1-(1-\beta)r\bigr]\Bigr).
  \label{eq:dCS-dPS-seller}
\end{equation}

%=============================================================================
\section{Pass-through}
%=============================================================================

At a zero rate, $\mu^{s}=\mu(p,h)$, so equilibrium satisfies
$D(p)=S(\mu(p,h))$. Write $Q\equiv D(p)$. A uniform one-dollar parallel shift
in every firm's marginal production cost, indexed by $c$, shifts market
supply at a given effective margin. A one-dollar increase in handling cost
$h$ lowers $\mu(p,h)$ at a fixed list price. The third shock is the ad
valorem fee at rate $\tau$.

For production and handling, market clearing is
\begin{equation}
  m(p,x)\equiv D(p)-S(\mu(p,x),x)=0,
  \qquad x\in\{c,h\}.
\end{equation}
Differentiating $m=0$ defines pass-through $\rho_x\equiv dp/dx$.
\begin{equation}
  \rho_x=-\frac{m_x}{m_p},
  \label{eq:pass-through-general}
\end{equation}
Expanding $m=D-S$ gives
\begin{align}
  m_p &= D'(p)-S'(\mu)\,\mu_p,
  &
  m_x &= -S'(\mu)\,\mu_x-S_x,
\end{align}
hence
\begin{equation}
  \rho_x
  = \frac{-S'(\mu)\,\mu_x-S_x}{S'(\mu)\,\mu_p-D'(p)}.
  \label{eq:pass-through}
\end{equation}

The two shocks differ in which margin moves at fixed $p$ or fixed $\mu$.
\begin{align}
  \mu_c &= 0,
  &
  S_c &= -S'(\mu),
  &
  \mu_h &= -r(p),
  &
  S_h &= 0.
\end{align}
The production term is $S_c=-S'(\mu)$. A one-dollar parallel rise in every
firm's marginal cost is equivalent, at fixed $\mu$, to evaluating market
supply at $\mu-1$. Handling lowers the effective margin at a fixed list
price. Only returned units incur $h$, so a one-dollar rise lowers that
margin by $r(p)$. Quantity supplied then falls by $r(p)\,S'(\mu)$, the
fraction $r(p)$ of the production-cost numerator.

At a zero rate, pass-through of the fee equals
$p\bigl[1-(1-\beta)r\bigr]$ times production-cost pass-through. The fee
leaves the buyer's payment and the refund at $p$, so demand remains $D(p)$
and the return rate remains $r(p)$. At a fixed list price, the fee lowers the
effective margin by $p\bigl[1-(1-\beta)r\bigr]$ dollars along an unchanged
$S(\mu)$. Quantity supplied then falls by that amount times $S'(\mu)$. The
same gap is a production-cost rise of $p\bigl[1-(1-\beta)r\bigr]$ dollars,
so $\rho_{\tau}=p\bigl[1-(1-\beta)r\bigr]\rho_c$ at $\tau=0$. With no
returns, $r=0$, so $\beta$ drops out, the factor equals $p$, and
pass-through is $p$ times unit-cost pass-through \citep{weyl2013pass}.
\begin{result}[Pass-through into list price]
\label{res:pass-through}
  Production-cost pass-through is
  \begin{equation}
    \rho_c
    = \frac{S'(\mu)}{S'(\mu)\,\mu_p-D'(p)}.
    \label{eq:rho-c}
  \end{equation}
  The remaining shocks connect to this object by
  \begin{align}
    \rho_h &= r(p)\,\rho_c,
    &
    \rho_{\tau} &= p\bigl[1-(1-\beta)r(p)\bigr]\,\rho_c
    \quad\text{at }\tau=0.
    \label{eq:rho-links}
  \end{align}
\end{result}

\subsection{Elasticity representation}

The same pass-through formulas can be written in list-price elasticities.
Write $\mu_p\equiv d\mu/dp$ for the response of the effective margin to the
list price. Demand and supply elasticities with respect to list price are
\[
\varepsilon_D(p)\equiv\frac{p D'(p)}{D(p)}<0,
\qquad
\varepsilon_S(p)\equiv\frac{dS(\mu(p))}{dp}\frac{p}{S(\mu(p))}
=S'(\mu)\,\mu_p\,\frac{p}{S(\mu)}.
\]
Here $\varepsilon_S(p)$ is the observed elasticity of market supply with
respect to the clearing price. It has the sign of $\mu_p$. With no returns,
$\mu_p=1$ and $\varepsilon_S>0$. With returns, $\mu_p$ can be negative. Then
$\varepsilon_S<0$, and a higher list price reduces the quantity firms supply.

With no returns, pass-through equals
$\varepsilon_S/(-\varepsilon_D+\varepsilon_S)$. Supply is then a function of
the list price, so $S'=dS/dp$. A one-dollar cost increase cuts supply by
$S'$ units at the original price and opens an excess-demand gap of that
size. A one-dollar rise in the list price closes the gap from both sides.
Demand falls by $|D'|$ and supply rises by $S'$. Closing a gap of size $S'$
therefore takes a price rise of $S'/(S'+|D'|)$.

Under the return margin, supply depends on the effective margin, so
$S'(\mu)=dS/d\mu$. A one-dollar production-cost increase leaves $\mu$
unchanged at a fixed list price. It is equivalent to evaluating supply at
$\mu-1$, so supply falls by $S'(\mu)$ units and opens an excess-demand gap
of that size. A one-dollar rise in the list price
closes the gap at rate $S'(\mu)\,\mu_p-D'(p)$, so
\[
\rho_c=\frac{S'(\mu)}{S'(\mu)\,\mu_p-D'(p)}.
\]
A larger $\mu_p$ means that one dollar of list price restores more of the
effective margin, and therefore more quantity supplied, so a smaller
list-price movement closes the gap. Writing the same formula with
$\varepsilon_S$ places $1/\mu_p$ in the numerator,
\begin{equation}
  \rho_c
  =
  \frac{\varepsilon_S(p)/\mu_p}
       {-\varepsilon_D(p)+\varepsilon_S(p)}.
  \label{eq:rho_c_elasticity_final}
\end{equation}

The sign of $\rho_c$ is the sign of the denominator
$-\varepsilon_D+\varepsilon_S$. The numerator $\varepsilon_S/\mu_p$ is
positive, because $\varepsilon_S$ and $\mu_p$ have the same sign.
\begin{result}[Negative pass-through]
\label{res:negative-pt}
  We have $\rho_c<0$ if and only if $\varepsilon_S<0$ ($\mu_p<0$) and
  $|\varepsilon_S|>-\varepsilon_D$.
\end{result}
A cost increase is then cleared by a lower list price. With
$\varepsilon_S<0$, so $\mu_p<0$, a higher list price lowers the quantity
firms supply. A rise in $c$ shifts that supply curve to the left and creates
excess demand at the initial $p$. A higher price would shrink supply further
and widen the gap. The price falls. A lower $p$ raises quantity demanded
and, because supply slopes down, raises quantity supplied. When
$|\varepsilon_S|>-\varepsilon_D$, supply is steeper than demand, so the
supply response to the price cut exceeds the demand response, and the excess
demand clears at a lower list price. In the classical model, $\mu_p=1$,
supply slopes up in the list price, and $\rho_c>0$.
\begin{proof}
By (\ref{eq:rho_c_elasticity_final}),
\[
  \rho_c = \frac{\varepsilon_S/\mu_p}{-\varepsilon_D+\varepsilon_S}.
\]
Since
$\varepsilon_S=S'(\mu)\mu_p\,p/S(\mu)$, the terms $\varepsilon_S$ and $\mu_p$
have the same sign, so $\varepsilon_S/\mu_p>0$.  Hence $\rho_c<0$ if and only
if
\[
  -\varepsilon_D+\varepsilon_S<0,
\]
which is equivalent to $\varepsilon_S<0$ and
$|\varepsilon_S|>-\varepsilon_D$.  Because $\varepsilon_S$ and $\mu_p$ have the
same sign, $\varepsilon_S<0$ is equivalent to $\mu_p<0$.
\end{proof}

The effective margin rises by more than one dollar per dollar of list price
when the return rate falls steeply enough. Differentiating
$\mu=p-r(p)(p+h)$ gives
\[
\mu_p=1-r(p)-r'(p)(p+h).
\]
Write $\varepsilon_r(p)\equiv p\,r'(p)/r(p)$ for the elasticity of the return
rate in list price. Then $\mu_p>1$ if and only if
$\varepsilon_r(p)<-p/(p+h)$. Selection must then outweigh the behavioral rise
in individual return probabilities by enough that the aggregate return rate
falls steeply in $p$. When the behavioral response dominates,
$\varepsilon_r$ is less negative, or positive, and $\mu_p<1$.

The classical no-return formula treats the observed list-price slope of
supply as if the effective margin moved one for one with price. That slope
is $S'(\mu)\,\mu_p$, so the formula returns $\mu_p\rho_c$.
\begin{result}
\label{res:pt-bias-er}
  Relative to the correct returns-adjusted $\rho_c$, the classical no-return
  pass-through formula overestimates pass-through when $\mu_p>1$, agrees
  with it when $\mu_p=1$, and underestimates it when $0\leq\mu_p<1$. When
  $\mu_p<0$, the classical formula has the opposite sign of $\rho_c$. The
  condition $\mu_p>1$ is equivalent to $\varepsilon_r(p)<-p/(p+h)$.
\end{result}
When $\mu_p>1$, one dollar of list price restores more than one dollar of
effective margin, so a smaller price movement closes the excess-demand gap
than the classical formula predicts. Higher handling cost widens the set of
$\varepsilon_r$ for which that overestimation occurs, because the threshold
$-p/(p+h)$ moves toward zero. When $0\leq\mu_p<1$, one dollar of list price
restores less than one dollar of effective margin, so the list price must
move by more than the classical formula predicts. When $\mu_p<0$, the
effective margin falls in the list price, and the classical formula predicts
a price change of the wrong sign.

%=============================================================================
\section{Local incidence}
%=============================================================================

Production cost, handling cost, and the fee reach consumers only
through the list price. Throughout, $Q\equiv D(p)$ and $r\equiv r(p)$ denote
equilibrium gross quantity and the return rate, and $\rho_x\equiv dp/dx$.
Consumer surplus therefore changes by
\begin{equation}
  \frac{dCS}{dx} = -Q(1-r)\rho_x,
  \label{eq:CSx_incidence}
\end{equation}

Producer surplus is $PS(\mu,x)=\int_0^{\mu} S(z,x)\,dz$. Differentiating in
$x$ and using the equilibrium condition $S(\mu,x)=Q$ gives
\[
  \frac{dPS}{dx}
  = Q\frac{d\mu}{dx} + \Sigma_x,
\]
where $\Sigma_x\equiv\int_0^\mu (\partial S(z,x)/\partial x)\,dz$ is the
direct shift in the area under the supply curve. The chain rule is
$\frac{d\mu}{dx} = \mu_p \rho_x + \mu_x$, so
\begin{equation}
  \frac{dPS}{dx} = Q(\mu_p \rho_x + \mu_x) + \Sigma_x.
  \label{eq:PSx_incidence}
\end{equation}

The local incidence ratio is $I_x \equiv (dCS/dx)/(dPS/dx)$. For these three
shocks,
\[
  I_x = \frac{-Q(1-r)\rho_x}{Q(\mu_p \rho_x + \mu_x) + \Sigma_x},
\]
and the rows of Table~\ref{tab:incidence_components} reduce to one expression
in $\rho_c$ and $\mu_p$.

\begin{table}[htbp]
\centering
\begin{tabular}{lccccccc}
\hline
Shock $x$ & $\rho_x$ & $\mu_x$ & $\Sigma_x$ & $dCS/dx$ & $dPS/dx$ & $I_x$ \\
\hline
Production cost $c$ & $\rho_c$ & $0$ & $-Q$ & $-Q(1-r)\rho_c$ & $Q(\mu_p\rho_c - 1)$ & $\displaystyle \frac{(1-r)\rho_c}{1-\mu_p\rho_c}$ \\
Handling cost $h$ & $r\rho_c$ & $-r$ & $0$ & $-Q(1-r)r\rho_c$ & $Qr(\mu_p\rho_c - 1)$ & $\displaystyle \frac{(1-r)\rho_c}{1-\mu_p\rho_c}$ \\
Fee $\tau$ & $\theta\rho_c$ & $-\theta$ & $0$ & $-Q(1-r)\theta\rho_c$ & $\theta Q(\mu_p\rho_c - 1)$ & $\displaystyle \frac{(1-r)\rho_c}{1-\mu_p\rho_c}$ \\
\hline
\end{tabular}
\caption{Pass-through, direct effects, surplus changes, and incidence ratio for
each shock.  The fee column is evaluated at a zero rate, with
$\theta\equiv p\bigl[1-(1-\beta)r\bigr]$.}
\label{tab:incidence_components}
\end{table}

\begin{result}
\label{res:incidence-ratio}
For production cost, handling cost, and the fee, the local incidence
ratio is identical:
\begin{equation}
  I_x = \frac{(1-r)\rho_c}{1-\mu_p \rho_c}>0.
  \label{eq:incidence_result}
\end{equation}
\end{result}

Handling cost and the fee shift the effective margin by a scalar, and that
scalar multiplies the changes in consumer and producer surplus. At a fixed
list price a production-cost shock leaves $\mu$ unchanged. It shifts supply,
and the producer loss before the price adjusts is $Q$, the same loss as a
one-dollar fall in the effective margin. Handling shifts $\mu$ down by $r$
dollars. At a zero rate the fee shifts it down by
$p\bigl[1-(1-\beta)r\bigr]$ dollars. For a small shock that direct shift
scales pass-through, so $\rho_h=r\rho_c$ and
$\rho_{\tau}=p\bigl[1-(1-\beta)r\bigr]\rho_c$ at $\tau=0$.

The scalar cancels in the incidence ratio. How far the equilibrium moves
depends on the shock. How the movement splits between consumers and
producers depends on the local slopes of demand and supply and on
$\mu_p$.\footnote{Consumers lose or gain through the list-price change
$\rho_x$. Producers' receipt moves with the effective margin, which changes
by $\mu_p\rho_x$ when the list price changes by $\rho_x$. Thus $\mu_p$
governs how much of each dollar of $p$ feeds into producer surplus.}

The denominator $1-\mu_p\rho_c$ in \eqref{eq:incidence_result} is the
per-unit producer loss on a one-dollar cost shock. Before the price
adjusts, producer surplus falls by $Q$. The list-price change moves the
effective margin by $\mu_p\rho_c$ per unit, and the envelope $dPS/d\mu=Q$
turns that movement into a recovery of $Q\mu_p\rho_c$. The net change in
producer surplus is $Q(\mu_p\rho_c-1)$, so the loss per unit is
$1-\mu_p\rho_c$. The numerator $(1-r)\rho_c$ is the per-unit consumer loss
on kept units. The incidence ratio compares these two per-unit burdens. When
the denominator is small and positive, consumers bear most of the loss.
When the denominator is negative, the recovery exceeds the direct cost and
producers gain.

Producer surplus rises with a production-cost increase exactly when
pass-through is negative. The recovery factor equals
$\mu_p\rho_c=\varepsilon_S/(\varepsilon_S-\varepsilon_D)$. It exceeds one if
and only if $\varepsilon_S<0$ and $|\varepsilon_S|>-\varepsilon_D$.
Result~\ref{res:negative-pt} states that this is the condition for
$\rho_c<0$.

When pass-through is negative, clearing requires the effective margin to
rise by more than one dollar. A one-dollar cost increase shifts the supply
curve upward and opens an excess-demand gap at the initial price. Quantity
supplied increases with the effective margin, so covering the original gap
requires $\mu$ to rise by one dollar. The price fall also raises quantity
demanded and opens a further gap. Supply must increase by more than the
original shift, so the effective margin must rise by more than one dollar.
Thus $\mu_p\rho_c>1$ is the market-clearing requirement when pass-through is
negative.

The incidence ratio is positive, so consumer and producer surplus move
together. When $\rho_c<0$, an increase in production cost, handling cost, or
the fee lowers the list price. Consumer surplus rises, and producer
surplus rises as well. When $\rho_c>0$, all three shocks reduce both
surpluses. In the classical model, $\mu_p=1$ and pass-through is positive,
so both surpluses fall.

\subsection{First-order private surplus and deadweight loss}
\label{sec:DWL}

Private surplus is the sum of consumer and producer surplus,
$W\equiv CS+PS$. Adding (\ref{eq:CSx_incidence}) and
(\ref{eq:PSx_incidence}) gives
\[
  \frac{dW}{dx}
  = -Q(1-r)\rho_x
    + Q(\mu_p\rho_x + \mu_x)
    + \Sigma_x
  = Q\bigl[\mu_x + \bigl(\mu_p-(1-r)\bigr)\rho_x\bigr]
    + \Sigma_x.
\]
Substituting $\mu_p=1-r-r'(p)(p+h)$ gives
\begin{equation}
  \frac{dW}{dx}
  = Q\bigl[\mu_x - r'(p)(p+h)\rho_x\bigr] + \Sigma_x,
  \label{eq:DWL_incidence}
\end{equation}
with $\mu_x$ and $\Sigma_x$ as in Table~\ref{tab:incidence_components}.
The term $-Q\,r'(p)(p+h)\rho_x$ is the first-order effect of the
price-induced change in the return rate. A change $r'(p)\rho_x$ in the
return rate costs $p+h$ per unit sold. The term is zero when $r'=0$.

With no returns, a unit production-cost increase lowers private surplus by
$Q$. In that benchmark $r=0$ and $r'=0$, so \eqref{eq:DWL_incidence} reduces
to
\[
  \frac{dW}{dx} = Q\,\mu_x + \Sigma_x.
\]
A unit production-cost increase shifts the supply curve up by one dollar at
every quantity. The private-surplus loss is a rectangle of height one and
width $Q$. Private surplus in that benchmark is $W=CS(p)+PS(p-c)$.
Differentiating, the price terms $-Q\rho_c$ and $+Q\rho_c$ cancel, so
$dW/dc=-Q$.

For an ad valorem tax at a zero rate and with no returns, $r=0$, so $\beta$
drops out. The same cancellation gives $dW/d\tau=-pQ$. One percentage point
is a wedge of $p$ dollars on every unit, and the price-change terms cancel.
Platform revenue $G=\tau p Q$ is a transfer. At $\tau=0$,
$dV/d\tau=dW/d\tau+pQ=0$. A one-dollar cost shock destroys $Q$ dollars of
private surplus. An infinitesimal fee moves $pQ$ dollars from private
surplus to the platform, so social surplus is unchanged at first order.

With returns, the return-rate term in \eqref{eq:DWL_incidence} is generally
nonzero for a production-cost increase, a handling-cost increase, and the
the fee. Private surplus then changes through the keep-or-return margin as
well as through the direct wedge. For the fee, platform revenue is still a
transfer. It is added back below, with
$G=\tau p\bigl[1-(1-\beta)r\bigr]Q$.

Substituting $\mu_x$ and $\Sigma_x$ from
Table~\ref{tab:incidence_components} into \eqref{eq:DWL_incidence}, and using
the pass-through links in Result~\ref{res:pass-through}, gives the following
derivatives. The tax derivative is evaluated at a zero rate.
\begin{align}
  \frac{dW}{dc}
  &= -Q\,\Lambda,
  &
  \frac{dW}{dh}
  &= -Q\,r\,\Lambda,
  &
  \frac{dW}{d\tau}
  &= -p\bigl[1-(1-\beta)r\bigr]Q\,\Lambda,
  \label{eq:DWL_shocks}
\end{align}
where
\begin{equation}
  \Lambda \equiv 1 + r'(p)(p+h)\,\rho_c.
  \label{eq:Lambda}
\end{equation}

The unit term in $\Lambda$ is the direct wedge on units already sold. For a
production-cost shock that wedge is one dollar. For the fee it is
$p\bigl[1-(1-\beta)r\bigr]$ dollars. List-price terms cancel, so the change
in the number of buyers does not enter at first order. The term
$r'(p)(p+h)\rho_c$ is the further change on the keep-or-return margin after
the price moves. Each change $r'(p)\rho_c$ in the return rate costs $p+h$
per unit sold. A handling-cost shock multiplies the production-cost
derivative by $r$.

Freeze the aggregate return rate at its initial value $r_0\equiv r(p_0)$.
In that ex-ante economy the effective margin is
\eqref{eq:mu-def} with $r(p)$ replaced by the constant $r_0$,
\begin{equation}
  \mu^{\mathrm{EA}}(p)\equiv p-r_0(p+h).
  \label{eq:mu-EA}
\end{equation}
Applying \eqref{eq:DWL_incidence} to the production-cost shock in this
economy yields $dW/dc\big|_{\mathrm{EA}}=-Q$, the direct cost on units sold
with the return rate held at $r_0$. The ex-post piece is the residual
$dW/dc\big|_{\mathrm{EP}}\equiv dW/dc-dW/dc\big|_{\mathrm{EA}}$. Since
$dW/dc=-Q\Lambda$,
\begin{equation}
  \frac{dW}{dc}\Big|_{\mathrm{EP}}
  =-Q\,r'(p)(p+h)\rho_c.
  \label{eq:DWL-EP}
\end{equation}
For the fee at a zero rate, both pieces are multiplied by
$p\bigl[1-(1-\beta)r\bigr]$. The ex-ante piece is
$-p\bigl[1-(1-\beta)r\bigr]Q$. The ex-post piece is
$-p\bigl[1-(1-\beta)r\bigr]Q\,r'(p)(p+h)\rho_c$.
\begin{result}[Ex-ante and ex-post decomposition]
\label{res:EA-EP}
  For a unit production-cost shock,
  \begin{equation}
    \frac{dW}{dc}
    =-Q
    -Q\,r'(p)(p+h)\rho_c.
    \label{eq:DWL-decomp}
  \end{equation}
  The first term is the ex-ante change on the purchase margin.  The second is
  the ex-post change on the keep-or-return margin.  Each unit switched from keep
  to return costs $p+h$.  For the fee at a zero rate both terms
  are multiplied by $p\bigl[1-(1-\beta)r\bigr]$.  For a handling shock the two terms become $-Qr$ and
  $-Qr\,r'(p)(p+h)\rho_c$.
  The ex-ante term is the direct cost on units already sold. The change in the
  number of buyers does not enter at first order, so this term is not a loss
  of purchases.
  The ex-post term is a further loss when $r'(p)>0$ and $\rho_c>0$: the higher
  price raises the return rate, and there are too many returns.
  It is a gain, with private surplus still falling, when $r'(p)<0$: pass-through
  is then positive, the higher price lowers the return rate, and there are fewer
  returns than at the original price.
  It is a gain that exceeds the ex-ante loss when $\rho_c<0$: the price fall
  cuts returns, and private surplus rises.
\end{result}

Consumer and producer surplus move together, so $\Lambda$ has the same sign
as $\rho_c$. The no-return formula keeps only the ex-ante rectangle $-Q$.
The gap between that formula and $dW/dc=-Q\Lambda$ is the ex-post term
$-Q\,r'(p)(p+h)\rho_c$. That term can offset the rectangle, amplify it, or
reverse the sign of the private-surplus change.

When selection dominates, $r'(p)<0$. Then $\mu_p>0$ and hence $\rho_c>0$, so
a cost increase raises the list price. The higher price screens out frequent
returners and lowers the return rate. The ex-post term is then a gain, and
$0<\Lambda<1$. The private-surplus change stays negative, because
$\rho_c>0$ keeps consumer and producer surplus falling.

When $r'(p)>0$ and $\rho_c>0$, the higher list price raises the return rate.
Extra returns cost $p+h$ each, so the ex-post term amplifies the rectangle
and $\Lambda>1$. That amplification is larger when $r'(p)$ or $h$ is large.

When $\rho_c<0$, a production-cost, handling-cost, or fee increase
lowers the list price, so consumer surplus rises. The incidence ratio is
positive, so producer surplus rises as well, a Pareto improvement. Negative
pass-through requires $r'(p)>0$, so the fall in the list price lowers the
return rate. Each unit that switches from return to keep saves $p+h$, and
that saving more than offsets the direct wedge. Private surplus rises by
$Q|\Lambda|$ under the cost shock and by
$p\bigl[1-(1-\beta)r\bigr]Q|\Lambda|$ under the fee, with $\Lambda<0$. The
gain is larger when $r'(p)(p+h)\,|\rho_c|$ is large.
Result~\ref{res:welfare-bias} records these three comparisons.
\begin{result}[No-return private-surplus bias]
\label{res:welfare-bias}
  Consider a unit production-cost increase.  Relative to the
  correct returns-adjusted change $dW/dc=-Q\Lambda$, the no-return formula
  $dW/dc=-Q$
  \begin{enumerate}
  \item[(i)] overestimates the private-surplus loss if and only if
  $r'(p)<0$ (which implies $\mu_p>0$ and hence $\rho_c>0$, so
  $0<\Lambda<1$);
  \item[(ii)] underestimates the private-surplus loss if and only if
  $r'(p)>0$ and $\rho_c>0$ (hence $\Lambda>1$);
  \item[(iii)] overestimates the private-surplus loss with the wrong
  sign---reporting a loss of $Q$ when the shock in fact raises private
  surplus by $Q|\Lambda|$---if and only if $r'(p)>0$ and $\rho_c<0$.
  \end{enumerate}
  For the fee at a zero rate the same three comparisons apply to the ex-ante
  benchmark $-p\bigl[1-(1-\beta)r\bigr]Q$, which holds the return rate fixed.
  A formula that also sets the return rate to zero uses $-pQ$. The two agree
  when $\beta=1$ or $r=0$. When $\beta<1$ and $r>0$, $-pQ$ is a larger private
  loss, because the ex-ante benchmark gives the seller back the share
  $1-\beta$ of the fee on each returned unit.
\end{result}

Social surplus adds back the revenue the platform keeps.
On a kept unit that revenue is $\tau p$. On a returned unit it is the share
$\beta\tau p$. Expected revenue per unit is therefore
$\tau p\bigl[1-(1-\beta)r\bigr]$, and
\[
  G=\tau p\bigl[1-(1-\beta)r\bigr]Q,
\]
with $Q$ the equilibrium gross quantity. Write $V=W+G$. Differentiating
at $\tau=0$ gives $dG/d\tau=p\bigl[1-(1-\beta)r\bigr]Q$, so
$dV/d\tau=dW/d\tau+p\bigl[1-(1-\beta)r\bigr]Q$. Using \eqref{eq:DWL_shocks} gives the change in social surplus.
\begin{result}[Social surplus]
\label{res:social}
At a zero rate,
\begin{equation}
  \frac{dV}{d\tau}
  = -p\bigl[1-(1-\beta)r\bigr]Q(\Lambda-1)
  = -p\bigl[1-(1-\beta)r\bigr]Q\,r'(p)(p+h)\rho_c.
  \label{eq:DWL_V_tax}
\end{equation}
An infinitesimal fee raises social surplus when $r'(p)\rho_c<0$ and lowers
it when $r'(p)\rho_c>0$. The factor $p\bigl[1-(1-\beta)r\bigr]$ scales this
change, so a larger share of the fee kept on a refund enlarges the gain or
the loss.
\end{result}

%=============================================================================
\section{Global Incidence}
\label{sec:global-incidence}
%=============================================================================

The fee row of Table~\ref{tab:incidence_components} is the derivative at a
zero rate. It is not the integrand for a finite change in $\tau$.

A finite production-cost or handling-cost change integrates that shock's row
in the same table, with the row re-evaluated at the current equilibrium.

Along such a path, the global incidence ratio
$I_{x_0}^{x_1}\equiv \Delta CS/\Delta PS$, when $\Delta PS\neq 0$, is the
weighted average of the local ratio implied by $dCS/dx = I(x)\,dPS/dx$,
\begin{equation}
I_{x_0}^{x_1} = \int_{x_0}^{x_1} I(x)\,\omega(x)\,dx,
\label{eq:global-incidence-final}
\end{equation}
where $\omega(x)\equiv (dPS/dx)/\Delta PS$.

The global ratio can be negative, even though $I(x)$ is positive at every
point on the path. If $dPS/dx$ changes sign along the path, the weights
$\omega(x)$ take both signs, and the weighted average of positive local
ratios can be negative. A finite production-cost or handling-cost change
therefore need not preserve the sign of local incidence.

\clearpage
\appendix
\section{Welfare bias in primitives}
\label{app:welfare-bias-primitives}
Result~\ref{res:welfare-bias} has an equivalent statement in elasticities.
Recall $\varepsilon_r(p)\equiv p\,r'(p)/r(p)$ and
$\mu_p=1-r(p)-r'(p)(p+h)$. Then
$r'(p)(p+h)=\varepsilon_r(p)\,r(p)\,(p+h)/p$, so
\begin{equation}
  \mu_p>0
    \;\iff\;
  \varepsilon_r(p)
  <
  \frac{p}{p+h}\,\frac{1-r(p)}{r(p)},
  \label{eq:app-mu-er}
\end{equation}

In these elasticities, Result~\ref{res:welfare-bias} reads as follows.
\begin{enumerate}
\item[(i)] overestimation iff $\varepsilon_r(p)<0$;
\item[(ii)] underestimation iff $\varepsilon_r(p)>0$ and either $\mu_p\geq 0$,
or both $\mu_p<0$ and $|\varepsilon_S|< -\varepsilon_D$;
\item[(iii)] overestimation with wrong sign (predicted loss, actual gain)
iff $\mu_p<0$ and $|\varepsilon_S|>-\varepsilon_D$.
\end{enumerate}
\begin{proof}
By Result~\ref{res:welfare-bias}, overestimation with $0<\Lambda<1$ is
$r'<0$, underestimation is $r'>0$ with $\rho_c>0$, and overestimation with
wrong sign is $r'>0$ with $\rho_c<0$.
Since $\varepsilon_r$ has the same sign as $r'$, write these in
$\varepsilon_r$.  For~(i), $\varepsilon_r<0$ forces $\mu_p>1-r>0$ and hence
$\rho_c>0$, so $0<\Lambda<1$.
When $\mu_p=0$ and $r(p)<1$, the derivative $r'(p)$ is positive,
$\varepsilon_S=0$, and $\rho_c=S'(\mu)/(-D'(p))>0$, so that case is in~(ii).
Result~\ref{res:negative-pt} translates the sign of $\rho_c$ into conditions
on $\mu_p$ and $|\varepsilon_S|\gtrless -\varepsilon_D$: with
$\varepsilon_r>0$ this is~(ii), and $\rho_c<0$ is~(iii).
\end{proof}

\bibliographystyle{plainnat}
\bibliography{references}

\end{document}
